%% Generated by Sphinx.
\def\sphinxdocclass{report}
\documentclass[letterpaper,10pt,english]{sphinxmanual}
\ifdefined\pdfpxdimen
   \let\sphinxpxdimen\pdfpxdimen\else\newdimen\sphinxpxdimen
\fi \sphinxpxdimen=.75bp\relax
\ifdefined\pdfimageresolution
    \pdfimageresolution= \numexpr \dimexpr1in\relax/\sphinxpxdimen\relax
\fi
%% let collapsible pdf bookmarks panel have high depth per default
\PassOptionsToPackage{bookmarksdepth=5}{hyperref}

\PassOptionsToPackage{booktabs}{sphinx}
\PassOptionsToPackage{colorrows}{sphinx}

\PassOptionsToPackage{warn}{textcomp}
\usepackage[utf8]{inputenc}
\ifdefined\DeclareUnicodeCharacter
% support both utf8 and utf8x syntaxes
  \ifdefined\DeclareUnicodeCharacterAsOptional
    \def\sphinxDUC#1{\DeclareUnicodeCharacter{"#1}}
  \else
    \let\sphinxDUC\DeclareUnicodeCharacter
  \fi
  \sphinxDUC{00A0}{\nobreakspace}
  \sphinxDUC{2500}{\sphinxunichar{2500}}
  \sphinxDUC{2502}{\sphinxunichar{2502}}
  \sphinxDUC{2514}{\sphinxunichar{2514}}
  \sphinxDUC{251C}{\sphinxunichar{251C}}
  \sphinxDUC{2572}{\textbackslash}
\fi
\usepackage{cmap}
\usepackage[T1]{fontenc}
\usepackage{amsmath,amssymb,amstext}
\usepackage{babel}



\usepackage{tgtermes}
\usepackage{tgheros}
\renewcommand{\ttdefault}{txtt}



\usepackage[Bjarne]{fncychap}
\usepackage{sphinx}

\fvset{fontsize=auto}
\usepackage{geometry}


% Include hyperref last.
\usepackage{hyperref}
% Fix anchor placement for figures with captions.
\usepackage{hypcap}% it must be loaded after hyperref.
% Set up styles of URL: it should be placed after hyperref.
\urlstyle{same}

\addto\captionsenglish{\renewcommand{\contentsname}{Contents:}}

\usepackage{sphinxmessages}
\setcounter{tocdepth}{1}



\title{ESP32\sphinxhyphen{}C3 Breadboard Adapter}
\date{Oct 19, 2024}
\release{5.3}
\author{Alexander Bobkov}
\newcommand{\sphinxlogo}{\vbox{}}
\renewcommand{\releasename}{Release}
\makeindex
\begin{document}

\ifdefined\shorthandoff
  \ifnum\catcode`\=\string=\active\shorthandoff{=}\fi
  \ifnum\catcode`\"=\active\shorthandoff{"}\fi
\fi

\pagestyle{empty}
\sphinxmaketitle
\pagestyle{plain}
\sphinxtableofcontents
\pagestyle{normal}
\phantomsection\label{\detokenize{index::doc}}


\sphinxAtStartPar
Add your content using \sphinxcode{\sphinxupquote{reStructuredText}} syntax. See the
\sphinxhref{https://www.sphinx-doc.org/en/master/usage/restructuredtext/index.html}{reStructuredText}
documentation for details.

\sphinxstepscope


\chapter{About}
\label{\detokenize{about:about}}\label{\detokenize{about::doc}}
\sphinxstepscope


\chapter{Schematic}
\label{\detokenize{schematic:schematic}}\label{\detokenize{schematic::doc}}
\sphinxstepscope


\chapter{Examples}
\label{\detokenize{examples:examples}}\label{\detokenize{examples::doc}}
\sphinxAtStartPar
MicroPython Blinky Example

\sphinxAtStartPar
import machine
from machine import Pin
import time, math

\sphinxAtStartPar
ONBOARD\_LED = 10        \# GPIO10, PIN 7
ONBOARD\_BTN = 3         \# GPIO3, 13

\sphinxAtStartPar
\# Configure on\sphinxhyphen{}board LED and push button
\# Stated GPIOs correspond to the wiring schematic
onboard\_button = Pin(ONBOARD\_BTN, Pin.IN, Pin.PULL\_UP)

\sphinxAtStartPar
led = machine.PWM(ONBOARD\_LED, freq=1000)
\begin{description}
\sphinxlineitem{def pulse(l, t):}\begin{description}
\sphinxlineitem{for i in range(20):}
\sphinxAtStartPar
l.duty(int(math.sin(i/10 * math.pi) * 500 + 500))
time.sleep\_ms(t)

\end{description}

\sphinxAtStartPar
l.duty(0)

\end{description}

\sphinxAtStartPar
\# Interrupt function to turn LED ON when on\sphinxhyphen{}board button is pressed
def button\_interrupt(pin):
\begin{quote}

\sphinxAtStartPar
pulse(led, 70)
\end{quote}

\sphinxAtStartPar
def main():
\begin{quote}

\sphinxAtStartPar
\# Assign interrupt to the on\sphinxhyphen{}board push button
onboard\_button.irq(trigger=Pin.IRQ\_FALLING, handler=button\_interrupt)
\end{quote}
\begin{description}
\sphinxlineitem{if \_\_name\_\_ == ‘\_\_main\_\_’:}
\sphinxAtStartPar
main()

\end{description}

\sphinxstepscope


\chapter{Bill of Materials}
\label{\detokenize{bom:bill-of-materials}}\label{\detokenize{bom::doc}}


\renewcommand{\indexname}{Index}
\printindex
\end{document}